% Level 06 - The routine for every change
% Question answered: what do I do every time I modify something, from picking the task to merging?
% Facts: report team rules (Jira key MWPMWP-n in commits and PR titles, Definition of Done in Section 7.4, peer review required).
% Proposals: branch naming and the exact CI checks; both live in .github/workflows/ci.yml and Makefile.
\documentclass[tikz,border=10pt]{standalone}
% =============================================================================
% MethaneWatch diagram style  (shared by every diagram under diagrams/)
% -----------------------------------------------------------------------------
% Why one shared file: every diagram must look like part of the same set, and a
% colour change (for example a layer re-coloured) should be one edit, not 33.
%
% Usage, at the top of any diagram:
%   \documentclass[tikz,border=10pt]{standalone}
%   % =============================================================================
% MethaneWatch diagram style  (shared by every diagram under diagrams/)
% -----------------------------------------------------------------------------
% Why one shared file: every diagram must look like part of the same set, and a
% colour change (for example a layer re-coloured) should be one edit, not 33.
%
% Usage, at the top of any diagram:
%   \documentclass[tikz,border=10pt]{standalone}
%   % =============================================================================
% MethaneWatch diagram style  (shared by every diagram under diagrams/)
% -----------------------------------------------------------------------------
% Why one shared file: every diagram must look like part of the same set, and a
% colour change (for example a layer re-coloured) should be one edit, not 33.
%
% Usage, at the top of any diagram:
%   \documentclass[tikz,border=10pt]{standalone}
%   \input{../_common/mwstyle.tex}
%   \begin{document} ... \end{document}
%
% Colour meaning is fixed across the whole set (see diagrams/README.md). Two colours plus grey:
%   blue = MethaneWatch's own parts (application, data, models, science, services)
%   red  = evidence, security, audit and hard limits (the things that must never be bypassed)
%   grey = people, cloud and infrastructure, external systems, notes
% The older names green, orange, purple, teal and gold are kept as aliases so that no diagram
% source needs to change; they now resolve to blue (green, orange, purple) or grey (teal, gold).
% =============================================================================
\usepackage[T1]{fontenc}
\usepackage[scaled=0.95]{helvet}
\renewcommand{\familydefault}{\sfdefault}
\hyphenpenalty=10000 \exhyphenpenalty=10000 % never hyphenate words inside diagrams
\usepackage{pifont}
\usepackage{amssymb}
\usepackage{array}
\usepackage{booktabs}
\usetikzlibrary{positioning,arrows.meta,shapes.geometric,shapes.multipart,%
  shapes.misc,fit,backgrounds,calc,decorations.pathreplacing,matrix,patterns}

% ---- palette ----------------------------------------------------------------
\definecolor{mwink}{HTML}{1F2933}
\definecolor{mwgrey}{HTML}{6B7785}
\definecolor{mwline}{HTML}{4A5560}
\definecolor{mwlight}{HTML}{F2F4F7}
\definecolor{mwblue}{HTML}{1F6FB2}   \definecolor{mwblueL}{HTML}{E3EFFA}
\definecolor{mwred}{HTML}{B5483F}    \definecolor{mwredL}{HTML}{F8E3E0}
% aliases (see note above)
\colorlet{mwgreen}{mwblue}   \colorlet{mwgreenL}{mwblueL}
\colorlet{mworange}{mwblue}  \colorlet{mworangeL}{mwblueL}
\colorlet{mwpurple}{mwblue}  \colorlet{mwpurpleL}{mwblueL}
\colorlet{mwteal}{mwgrey}    \colorlet{mwtealL}{mwlight}
\colorlet{mwgold}{mwgrey}    \colorlet{mwgoldL}{mwlight}

% ---- node and edge styles ---------------------------------------------------
\tikzset{
  >={Stealth[length=2.2mm]},
  every node/.style={font=\sffamily\small,text=mwink},
  box/.style={rectangle,rounded corners=2.5pt,draw=mwline,line width=0.7pt,
    fill=white,align=center,inner sep=5pt,minimum height=7mm},
  blue/.style={draw=mwblue,fill=mwblueL},
  green/.style={draw=mwgreen,fill=mwgreenL},
  orange/.style={draw=mworange,fill=mworangeL},
  purple/.style={draw=mwpurple,fill=mwpurpleL},
  red/.style={draw=mwred,fill=mwredL},
  teal/.style={draw=mwteal,fill=mwtealL},
  gold/.style={draw=mwgold,fill=mwgoldL},
  grey/.style={draw=mwgrey,fill=mwlight},
  strong/.style={line width=1.6pt},
  ext/.style={box,dashed,draw=mwgrey,fill=mwlight},
  actor/.style={box,gold,rounded corners=8pt,minimum height=9mm},
  store/.style={cylinder,shape border rotate=90,aspect=0.22,draw=mwgreen,fill=mwgreenL,
    line width=0.7pt,align=center,inner sep=3pt,minimum height=11mm,minimum width=20mm},
  group/.style={rectangle,rounded corners=4pt,draw=mwline,dashed,line width=0.7pt,
    inner sep=7pt},
  glabel/.style={anchor=north west,font=\sffamily\bfseries\footnotesize,text=mwgrey,inner sep=2pt},
  note/.style={font=\sffamily\footnotesize,text=mwgrey,align=left,inner sep=1pt},
  tiny/.style={font=\sffamily\scriptsize},
  head/.style={font=\sffamily\bfseries\small},
  flow/.style={->,line width=0.8pt,draw=mwline},
  dflow/.style={->,line width=0.8pt,draw=mwline,dashed},
  bflow/.style={<->,line width=0.8pt,draw=mwline},
  lab/.style={font=\sffamily\scriptsize,text=mwgrey,fill=white,inner sep=1pt},
}

% ---- title / footer helpers -------------------------------------------------
% \dtitle{Finding or name}{One-line subtitle: what the diagram answers}
\newcommand{\dtitle}[2]{%
  \node[anchor=south west,font=\sffamily\bfseries\Large,text=mwink,inner sep=0pt] (dt)
    at ([yshift=12mm]current bounding box.north west) {#1};
  \node[anchor=north west,font=\sffamily\footnotesize,text=mwgrey,inner sep=0pt]
    at ([yshift=-1.2mm]dt.south west) {#2};}
% \dfoot{Source or reading note}
\newcommand{\dfoot}[2][14cm]{%
  \node[anchor=north west,note,text width=#1] at
    ([yshift=-4mm]current bounding box.south west) {#2};}
% in-node text helpers: \nm{Name} bold name, \tech{technology} small grey italic
\newcommand{\nm}[1]{{\bfseries\small #1}}
\newcommand{\tech}[1]{{\scriptsize\itshape\color{mwgrey} #1}}
% check mark used in matrices
\newcommand{\yes}{\ding{51}}

%   \begin{document} ... \end{document}
%
% Colour meaning is fixed across the whole set (see diagrams/README.md). Two colours plus grey:
%   blue = MethaneWatch's own parts (application, data, models, science, services)
%   red  = evidence, security, audit and hard limits (the things that must never be bypassed)
%   grey = people, cloud and infrastructure, external systems, notes
% The older names green, orange, purple, teal and gold are kept as aliases so that no diagram
% source needs to change; they now resolve to blue (green, orange, purple) or grey (teal, gold).
% =============================================================================
\usepackage[T1]{fontenc}
\usepackage[scaled=0.95]{helvet}
\renewcommand{\familydefault}{\sfdefault}
\hyphenpenalty=10000 \exhyphenpenalty=10000 % never hyphenate words inside diagrams
\usepackage{pifont}
\usepackage{amssymb}
\usepackage{array}
\usepackage{booktabs}
\usetikzlibrary{positioning,arrows.meta,shapes.geometric,shapes.multipart,%
  shapes.misc,fit,backgrounds,calc,decorations.pathreplacing,matrix,patterns}

% ---- palette ----------------------------------------------------------------
\definecolor{mwink}{HTML}{1F2933}
\definecolor{mwgrey}{HTML}{6B7785}
\definecolor{mwline}{HTML}{4A5560}
\definecolor{mwlight}{HTML}{F2F4F7}
\definecolor{mwblue}{HTML}{1F6FB2}   \definecolor{mwblueL}{HTML}{E3EFFA}
\definecolor{mwred}{HTML}{B5483F}    \definecolor{mwredL}{HTML}{F8E3E0}
% aliases (see note above)
\colorlet{mwgreen}{mwblue}   \colorlet{mwgreenL}{mwblueL}
\colorlet{mworange}{mwblue}  \colorlet{mworangeL}{mwblueL}
\colorlet{mwpurple}{mwblue}  \colorlet{mwpurpleL}{mwblueL}
\colorlet{mwteal}{mwgrey}    \colorlet{mwtealL}{mwlight}
\colorlet{mwgold}{mwgrey}    \colorlet{mwgoldL}{mwlight}

% ---- node and edge styles ---------------------------------------------------
\tikzset{
  >={Stealth[length=2.2mm]},
  every node/.style={font=\sffamily\small,text=mwink},
  box/.style={rectangle,rounded corners=2.5pt,draw=mwline,line width=0.7pt,
    fill=white,align=center,inner sep=5pt,minimum height=7mm},
  blue/.style={draw=mwblue,fill=mwblueL},
  green/.style={draw=mwgreen,fill=mwgreenL},
  orange/.style={draw=mworange,fill=mworangeL},
  purple/.style={draw=mwpurple,fill=mwpurpleL},
  red/.style={draw=mwred,fill=mwredL},
  teal/.style={draw=mwteal,fill=mwtealL},
  gold/.style={draw=mwgold,fill=mwgoldL},
  grey/.style={draw=mwgrey,fill=mwlight},
  strong/.style={line width=1.6pt},
  ext/.style={box,dashed,draw=mwgrey,fill=mwlight},
  actor/.style={box,gold,rounded corners=8pt,minimum height=9mm},
  store/.style={cylinder,shape border rotate=90,aspect=0.22,draw=mwgreen,fill=mwgreenL,
    line width=0.7pt,align=center,inner sep=3pt,minimum height=11mm,minimum width=20mm},
  group/.style={rectangle,rounded corners=4pt,draw=mwline,dashed,line width=0.7pt,
    inner sep=7pt},
  glabel/.style={anchor=north west,font=\sffamily\bfseries\footnotesize,text=mwgrey,inner sep=2pt},
  note/.style={font=\sffamily\footnotesize,text=mwgrey,align=left,inner sep=1pt},
  tiny/.style={font=\sffamily\scriptsize},
  head/.style={font=\sffamily\bfseries\small},
  flow/.style={->,line width=0.8pt,draw=mwline},
  dflow/.style={->,line width=0.8pt,draw=mwline,dashed},
  bflow/.style={<->,line width=0.8pt,draw=mwline},
  lab/.style={font=\sffamily\scriptsize,text=mwgrey,fill=white,inner sep=1pt},
}

% ---- title / footer helpers -------------------------------------------------
% \dtitle{Finding or name}{One-line subtitle: what the diagram answers}
\newcommand{\dtitle}[2]{%
  \node[anchor=south west,font=\sffamily\bfseries\Large,text=mwink,inner sep=0pt] (dt)
    at ([yshift=12mm]current bounding box.north west) {#1};
  \node[anchor=north west,font=\sffamily\footnotesize,text=mwgrey,inner sep=0pt]
    at ([yshift=-1.2mm]dt.south west) {#2};}
% \dfoot{Source or reading note}
\newcommand{\dfoot}[2][14cm]{%
  \node[anchor=north west,note,text width=#1] at
    ([yshift=-4mm]current bounding box.south west) {#2};}
% in-node text helpers: \nm{Name} bold name, \tech{technology} small grey italic
\newcommand{\nm}[1]{{\bfseries\small #1}}
\newcommand{\tech}[1]{{\scriptsize\itshape\color{mwgrey} #1}}
% check mark used in matrices
\newcommand{\yes}{\ding{51}}

%   \begin{document} ... \end{document}
%
% Colour meaning is fixed across the whole set (see diagrams/README.md). Two colours plus grey:
%   blue = MethaneWatch's own parts (application, data, models, science, services)
%   red  = evidence, security, audit and hard limits (the things that must never be bypassed)
%   grey = people, cloud and infrastructure, external systems, notes
% The older names green, orange, purple, teal and gold are kept as aliases so that no diagram
% source needs to change; they now resolve to blue (green, orange, purple) or grey (teal, gold).
% =============================================================================
\usepackage[T1]{fontenc}
\usepackage[scaled=0.95]{helvet}
\renewcommand{\familydefault}{\sfdefault}
\hyphenpenalty=10000 \exhyphenpenalty=10000 % never hyphenate words inside diagrams
\usepackage{pifont}
\usepackage{amssymb}
\usepackage{array}
\usepackage{booktabs}
\usetikzlibrary{positioning,arrows.meta,shapes.geometric,shapes.multipart,%
  shapes.misc,fit,backgrounds,calc,decorations.pathreplacing,matrix,patterns}

% ---- palette ----------------------------------------------------------------
\definecolor{mwink}{HTML}{1F2933}
\definecolor{mwgrey}{HTML}{6B7785}
\definecolor{mwline}{HTML}{4A5560}
\definecolor{mwlight}{HTML}{F2F4F7}
\definecolor{mwblue}{HTML}{1F6FB2}   \definecolor{mwblueL}{HTML}{E3EFFA}
\definecolor{mwred}{HTML}{B5483F}    \definecolor{mwredL}{HTML}{F8E3E0}
% aliases (see note above)
\colorlet{mwgreen}{mwblue}   \colorlet{mwgreenL}{mwblueL}
\colorlet{mworange}{mwblue}  \colorlet{mworangeL}{mwblueL}
\colorlet{mwpurple}{mwblue}  \colorlet{mwpurpleL}{mwblueL}
\colorlet{mwteal}{mwgrey}    \colorlet{mwtealL}{mwlight}
\colorlet{mwgold}{mwgrey}    \colorlet{mwgoldL}{mwlight}

% ---- node and edge styles ---------------------------------------------------
\tikzset{
  >={Stealth[length=2.2mm]},
  every node/.style={font=\sffamily\small,text=mwink},
  box/.style={rectangle,rounded corners=2.5pt,draw=mwline,line width=0.7pt,
    fill=white,align=center,inner sep=5pt,minimum height=7mm},
  blue/.style={draw=mwblue,fill=mwblueL},
  green/.style={draw=mwgreen,fill=mwgreenL},
  orange/.style={draw=mworange,fill=mworangeL},
  purple/.style={draw=mwpurple,fill=mwpurpleL},
  red/.style={draw=mwred,fill=mwredL},
  teal/.style={draw=mwteal,fill=mwtealL},
  gold/.style={draw=mwgold,fill=mwgoldL},
  grey/.style={draw=mwgrey,fill=mwlight},
  strong/.style={line width=1.6pt},
  ext/.style={box,dashed,draw=mwgrey,fill=mwlight},
  actor/.style={box,gold,rounded corners=8pt,minimum height=9mm},
  store/.style={cylinder,shape border rotate=90,aspect=0.22,draw=mwgreen,fill=mwgreenL,
    line width=0.7pt,align=center,inner sep=3pt,minimum height=11mm,minimum width=20mm},
  group/.style={rectangle,rounded corners=4pt,draw=mwline,dashed,line width=0.7pt,
    inner sep=7pt},
  glabel/.style={anchor=north west,font=\sffamily\bfseries\footnotesize,text=mwgrey,inner sep=2pt},
  note/.style={font=\sffamily\footnotesize,text=mwgrey,align=left,inner sep=1pt},
  tiny/.style={font=\sffamily\scriptsize},
  head/.style={font=\sffamily\bfseries\small},
  flow/.style={->,line width=0.8pt,draw=mwline},
  dflow/.style={->,line width=0.8pt,draw=mwline,dashed},
  bflow/.style={<->,line width=0.8pt,draw=mwline},
  lab/.style={font=\sffamily\scriptsize,text=mwgrey,fill=white,inner sep=1pt},
}

% ---- title / footer helpers -------------------------------------------------
% \dtitle{Finding or name}{One-line subtitle: what the diagram answers}
\newcommand{\dtitle}[2]{%
  \node[anchor=south west,font=\sffamily\bfseries\Large,text=mwink,inner sep=0pt] (dt)
    at ([yshift=12mm]current bounding box.north west) {#1};
  \node[anchor=north west,font=\sffamily\footnotesize,text=mwgrey,inner sep=0pt]
    at ([yshift=-1.2mm]dt.south west) {#2};}
% \dfoot{Source or reading note}
\newcommand{\dfoot}[2][14cm]{%
  \node[anchor=north west,note,text width=#1] at
    ([yshift=-4mm]current bounding box.south west) {#2};}
% in-node text helpers: \nm{Name} bold name, \tech{technology} small grey italic
\newcommand{\nm}[1]{{\bfseries\small #1}}
\newcommand{\tech}[1]{{\scriptsize\itshape\color{mwgrey} #1}}
% check mark used in matrices
\newcommand{\yes}{\ding{51}}

\begin{document}
\begin{tikzpicture}[
  s/.style={box,blue,text width=3.0cm,align=left,inner sep=4pt,minimum height=2.3cm,font=\sffamily\footnotesize},
  r/.style={box,grey,text width=3.0cm,align=left,inner sep=4pt,minimum height=2.3cm,font=\sffamily\footnotesize}]

  \node[glabel,anchor=south west] at (-0.2,1.5) {On your laptop};
  \node[s] (a) at (1.5,0)  {\nm{1  Pick the task}\\ Take the Jira key MWPMWP-n. Read the header of the file and its diagram.};
  \node[s] (b) at (5.2,0)  {\nm{2  Branch}\\ One branch per key, named with it.};
  \node[s] (c) at (8.9,0)  {\nm{3  Change}\\ Stay inside the header's import rules. Write or update the test named in docs/requirements.md.};
  \node[s] (d) at (12.6,0) {\nm{4  Check locally}\\ \texttt{make lint}, \texttt{make contracts}, \texttt{make test}. Model code: add \texttt{make test-ml}.};
  \node[s] (e) at (16.3,0) {\nm{5  Commit}\\ pre-commit runs ruff and refuses a message without MWPMWP-n.};
  \node[s] (f) at (20.0,0) {\nm{6  Open a pull request}\\ Key in the title; tick the Definition of Done list.};

  \node[glabel,anchor=south west] at (-0.2,-3.4) {On GitHub};
  \node[r] (g) at (20.0,-5.0) {\nm{7  CI}\\ ruff, mypy, import contracts, tests, web build, diagram build.};
  \node[r] (h) at (16.3,-5.0) {\nm{8  Peer review}\\ At least one teammate who is not the author.};
  \node[r,strong] (i) at (12.6,-5.0) {\nm{9  Merge to main}\\ Only when CI is green and the review is approved.};
  \node[r] (j) at (8.9,-5.0) {\nm{10  After merge}\\ CI builds the images; Beta only: deploy (see environments-promotion).};

  \foreach \x/\y in {a/b,b/c,c/d,d/e,e/f} {\draw[flow] (\x) -- (\y);}
  \draw[flow] (f.south) -- (g.north);
  \draw[flow] (g) -- (h); \draw[flow] (h) -- (i); \draw[flow] (i) -- (j);
  % loop back
  \draw[dflow,mwred] (g.south) -- ++(0,-0.9) -| node[lab,pos=0.12,above,text=mwred] {CI red or changes requested: back to step 3} ([xshift=-1.0cm]c.south);

  % extra duties
  \node[box,red,anchor=north west,text width=21.8cm,align=left,inner sep=6pt,font=\sffamily\footnotesize] at (-0.3,-8.1) {%
    \nm{Also do this when your change touches one of these}\\[3pt]
    \begin{tabular}{@{}p{5.2cm}p{16.2cm}@{}}
    A model or the trainer & Log the run in MLflow. Promotion is a separate act, only through the release gate (Beta: F1 $\ge$ 0.80, FPR $<$ 0.03\%).\\[2pt]
    Data, recipes or splits & Run \texttt{dvc repro}, commit the new manifest hash, update the dataset card. Never put synthetic ids in the benchmark.\\[2pt]
    The API contract (schemas, routes) & Run \texttt{scripts/export\_openapi.py} then \texttt{scripts/gen\_ts\_client.sh}; commit both outputs.\\[2pt]
    \texttt{configs/permissions.yaml} or \texttt{release\_gates.yaml} & Update the test that names the changed cell or number; say why in the pull request. Needs the owner of security (Ha) as reviewer.\\[2pt]
    A status, estimate or evidence rule & Change the rule test first. Never relax it to make a change pass.\\[2pt]
    A new package, port or route & Update the matching diagram and run \texttt{make diagrams}; add an ADR if it changes a decision.\\
    \end{tabular}};

  \dtitle{Every change takes the same ten steps}
         {Level 06 - change workflow. Blue: your laptop. Grey: GitHub. Red: what to add when the change touches a rule that must not weaken.}
  \dfoot[22cm]{Branch naming and the exact CI checks are proposals. The Jira key, the Definition of Done and peer review are team rules from the report.}
\end{tikzpicture}
\end{document}
